\documentclass[12pt]{article}
\usepackage[utf8]{inputenc}
\usepackage[T1]{fontenc}
\usepackage{hyperref}
\title{Ikpeng Language Resources for K-12 Curriculum Development}
\author{Compiled by Claude Sonnet 5 (claude-sonnet-5) for Andr\'as Kornai}
\date{}
\begin{document}
\maketitle

\section{Introduction}
Ikpeng (ISO 639-3 txi) is a Cariban language of the Xingu Indigenous Park in Brazil. Speaker estimates range from about 477 to 584 in a 2024 language snapshot, which describes the language as vital. Ikpeng teachers developed a Latin-script orthography in 1994.

\section{Primary Educational Resources}
\begin{itemize}
\item Ikpeng schools: three are reported, the largest being Escola Ind\'igena Central Ikpeng ``Amure'' in Moygu. The Ikpeng School is said to produce teaching materials for other Ikpeng communities. See \url{https://en.wikipedia.org/wiki/Ikpeng_language} and \url{https://www.omniglot.com/writing/ikpeng.htm}.
\end{itemize}
I found no titles of individual textbooks or primers.

\section{Linguistic and Reference Materials}
\begin{itemize}
\item Language Snapshot on Ikpeng, \emph{Language Documentation \& Description}: \url{https://www.lddjournal.org/article/id/1007}. It describes a 2009--2012 PRODOCLIN project (Museu do \'Indio/FUNAI with UNESCO). The project produced a sociolinguistic diagnosis, a lexical database, a descriptive grammar and a monolingual book of narratives. The authors note that open archiving has been hindered.
\item Article in \emph{Cadernos de Etnolingu\'istica / Abralin} (article 703): \url{https://cadernos.abralin.org/index.php/cadernos/article/view/703}
\item Omniglot on the Ikpeng writing system: \url{https://www.omniglot.com/writing/ikpeng.htm}
\item Ikpeng people (Wikipedia): \url{https://en.wikipedia.org/wiki/Ikpeng}
\end{itemize}
I found no published dictionary online.

\section{Religious and Community Resources}
I found no religious resources in my searches.

\section{Digital and Audio Resources}
I found no open audio or software resources. The lexical database mentioned above is not confirmed to be publicly available.

\section{Curriculum Development Context}
Ikpeng is still transmitted to children and has community schools, so a native-language curriculum is realistic. The main task is to obtain the existing PRODOCLIN outputs, the school-produced materials and the lexical database, ideally in cooperation with the Ikpeng teachers.
\end{document}